\section{Introduction}
Computational notebooks such as Jupyter have become central tools in scientific computing, teaching, and data analysis. Their interactive structure—mixing narrative text, code, and visual output—makes them highly suitable for inquiry-based learning in programming education~\cite{perkel2018jupyter}. Yet despite their pedagogical advantages, instructors often lack visibility into how students actually work within these environments. Students, on the other hand, often have little opportunity to reflect on their own workflow, error patterns, or problem-solving behavior. Research in programming education consistently shows that learners repeat similar types of mistakes, struggle to identify the sources of their errors, and often adopt inefficient trial-and-error strategies. At the same time, teachers typically only see final code submissions, which provide limited insight into the learners' process~\cite{robinson2022error}. This creates a gap between the correctness of the final code and the path the student took to reach it. Understanding this process is crucial to improve teaching interventions and help students develop their metacognitive skills.

Learning dashboards have been shown to enhance students' self-awareness, foster reflection, and offer instructors actionable insights into learner behavior. However, the majority of existing dashboards focus on high-level learning management system data rather than fine grained programming behavior within computational notebooks~\cite{paulsen2024learning}. 

This project addresses this gap by designing a JupyterLab dashboard extension capable of tracking detailed notebook interactions. The focus will be on the student dashboard, supporting self-reflection and awareness of one's coding patterns, errors, and work strategies. The overarching goal is to determine whether such dashboards can improve reflective learning, reduce repeated mistakes, and provide insights into students' perceptions of the usefulness and usability of the dashboard. To evaluate the effectiveness of this system, a controlled study will be conducted. Participants will complete two programming units in Jupyter Notebook. The first unit introduces basic Python concepts, while the second unit builds on these skills with slightly more advanced exercises. The experimental group will have access to the student dashboard for reflection during the learning units, and after each session, they will receive a more detailed view, while the control group will complete the tasks without feedback. Both quantitative measures, such as error reduction, execution efficiency, and task correctness, and qualitative measures, including self-reported reflection and usability feedback via a survey will be collected.

% Introduction
% - Significance
% - Motivation
% - Summary of Study Design